% Preface

\newglossaryentry{LHC}{
    name=LHC,
    description={Large Hadron Collider, a high-energy particle accelerator used for exploring fundamental interactions in physics},
    first={Large Hadron Collider (LHC)},
    text={LHC}
}

\newglossaryentry{ATLAS}{
    name=ATLAS,
    description={One of the main experiments at the LHC},
    first={ATLAS},
    text={ATLAS}
}

\newglossaryentry{pp}{
    name={$pp$},
    description={Proton-proton, Refers to proton-proton collisions, the primary collision type at the LHC},
    first={proton-proton ($pp$)},
    text={$pp$}
}

\newglossaryentry{GN++}{
    name=GN++,
    description={A custom graph neural network, with dedicated edge, node, and global updates, built for flavour tagging in ATLAS},
    first={GN++},
    text={GN++}
}

\newglossaryentry{GN1}{
    name=GN1,
    description={A graph neural network based flavour tagger released by ATLAS in 2022},
    first={GN1},
    text={GN1}
}

\newglossaryentry{GN2}{
    name=GN2,
    description={A transformer based flavour tagger, developed in this thesis, which is currently used by the ATLAS Collaboration},
    first={GN2},
    text={GN2}
}

\newglossaryentry{Spice}{
    name=Spice,
    description={A transformer model for flavour tagging},
    first={Spice},
    text={Spice}
}

\newglossaryentry{Salt}{
    name=\texttt{Salt},
    description={A repository for the continual development for transformer based models in ATLAS},
    first={\texttt{Salt}},
    text={\texttt{Salt}}
}

\newglossaryentry{PCJedi}{
    name=PCJedi,
    description={Particle Cloud Jet Diffusion, a deep generative model for high-energy physics jets},
    first={Point Cloud Jet Diffusion (\pcjedi)},
    text={\pcjedi}
}

\newglossaryentry{PCDroid}{
    name=PCDroid,
    description={Particle Cloud Diffusion of reconstructed objects with improved denoising, a deep generative model for high-energy physics jets and the successor to the PCJedi model},
    first={\pcdroid},
    text={\pcdroid}
}

\newglossaryentry{MPM}{
    name=MPM,
    description={Masked Particle Modelling, a pre-training strategy for high energy physics data based on a denoising autoencoder},
    first={Masked Particle Modelling (MPM)},
    text={MPM}
}

\newglossaryentry{MPM2}{
    name=MPM2,
    description={Masked Particle Modelling version-2, an improved version of the MPM pre-training strategy},
    first={Masked Particle Modelling Version-2 (MPM2)},
    text={MPM2}
}

% Standard Model of Particle Physics

\newglossaryentry{SM}{
    name=SM,
    description=The Standard Model of particle physics,
    first={Standard Model (SM)},
    text={SM}
}

\newglossaryentry{BSM}{
    name=BSM,
    description={Beyond the Standard Model, refers to theories and models that extend or go beyond the Standard Model of particle physics},
    first={beyond the Standard Model (BSM)},
    text={BSM}
}

\newglossaryentry{QFT}{
    name=QFT,
    description={Quantum Field Theory},
    first={Quantum Field Theory (QFT)},
    text={QFT}
}

\newglossaryentry{GSW}{
    name=GSW,
    description={The Glashow-Weinberg-Salam model, developed to unify the electromagnetic and weak forces},
    first={Glashow-Weinberg-Salam model (GSW)},
    text={GSW}
}

\newglossaryentry{QED}{
    name=QED,
    description={Quantum Electrodynamics, a complete description of the electromagnetic force},
    first={Quantum Electrodynamics (QED)},
    text={QED}
}

\newglossaryentry{QCD}{
    name=QCD,
    description={Quantum Chromodynamics, a renormalizable QFT describing quarks, gluons, and their interactions via the strong force},
    first={Quantum Chromodynamics (QCD)},
    text={QCD}
}

\newglossaryentry{CKM}{
    name=CKM,
    description={Cabibbo-Kobayashi-Maskawa matrix, describes quark mixing between weak and mass eigenstates},
    first={Cabibbo-Kobayashi-Maskawa (CKM)},
    text={CKM}
}

\newglossaryentry{PMNS}{
    name=PMNS,
    description={Pontecorvo-Maki-Nakagawa-Sakata matrix, describes lepton mixing},
    first={Pontecorvo-Maki-Nakagawa-Sakata (PMNS)},
    text={PMNS}
}

\newglossaryentry{EWSB}{
    name=EWSB,
    description={Electroweak Symmetry Breaking},
    first={Electroweak Symmetry Breaking (EWSB)},
    text={EWSB}
}

\newglossaryentry{CERN}{
    name=CERN,
    description={The European Organization for Nuclear Research},
    first={European Organization for Nuclear Research (CERN)},
    text={CERN}
}

\newglossaryentry{CMS}{
    name=CMS,
    description={Compact Muon Solenoid, a collider experiment at the LHC},
    first={CMS},
    text={CMS}
}

\newglossaryentry{GUTs}{
    name=GUTs,
    description={Grand Unified Theories, predict that all forces in the Standard Model are unified at high energies},
    first={Grand Unified Theories (GUTs)},
    text={GUTs}
}

% LHC and ATLAS

\newglossaryentry{LHCb}{
    name=LHCb,
    description={A collider experiment at the LHC focusing on measuring the properties of B mesons and other particles containing b-quarks},
    first={LHCb},
    text={LHCb}
}

\newglossaryentry{ALICE}{
    name=ALICE,
    description={A collider experiment at the LHC investigating the properties of quark-gluon plasma},
    first={ALICE},
    text={ALICE}
}

\newglossaryentry{IPs}{
    name=IPs,
    description={Interaction points within the LHC where particle collisions occur},
    first={interaction points (IPs)},
    text={IPs}
}

\newglossaryentry{Linac}{
    name=Linac,
    description={Linear Accelerator, a stage in the CERN accelerator complex that accelerates particles from rest},
    first={Linear Accelerator (Linac)},
    text={Linac}
}

\newglossaryentry{PSB}{
    name=PSB,
    description={Proton Synchrotron Booster, a stage in the CERN accelerator complex that increases particle energy},
    first={Proton Synchrotron Booster (PSB)},
    text={PSB}
}

\newglossaryentry{PS}{
    name=PS,
    description={Proton Synchrotron, a stage in the CERN accelerator complex that accelerates particles to higher energies},
    first={Proton Synchrotron (PS)},
    text={PS}
}

\newglossaryentry{SPS}{
    name=SPS,
    description={Super Proton Synchrotron, a stage in the CERN accelerator complex that accelerates particles before injection into the LHC},
    first={Super Proton Synchrotron (SPS)},
    text={SPS}
}

\newglossaryentry{ID}{
    name=ID,
    description={The innermost sub-detector of ATLAS, comprising the Pixel Detector, SCT, and TRT, used for tracking charged particles},
    first={Inner Detector (ID)},
    text={ID}
}

\newglossaryentry{SCT}{
    name=SCT,
    description={Semi-Conductor Tracker, a sub-detector within the ATLAS Inner Detector composed of silicon microstrip layers},
    first={Semi-Conductor Tracker (SCT)},
    text={SCT}
}

\newglossaryentry{TRT}{
    name=TRT,
    description={Transition Radiation Tracker, the outermost sub-detector of the ATLAS Inner Detector, used for tracking and particle identification},
    first={Transition Radiation Tracker (TRT)},
    text={TRT}
}

\newglossaryentry{IBL}{
    name=IBL,
    description={Insertable B-Layer, an additional pixel layer added to the ATLAS Inner Detector between Run 1 and Run 2},
    first={Insertable B-Layer (IBL)},
    text={IBL}
}

\newglossaryentry{ECal}{
    name=ECal,
    description={Electromagnetic Calorimeter, an ATLAS calorimeter system specialized for measuring the energy of electrons and photons},
    first={Electromagnetic Calorimeter (ECal)},
    text={ECal}
}

\newglossaryentry{EMEC}{
    name=EMEC,
    description={Liquid Argon Electromagnetic End-Cap, part of the ATLAS Electromagnetic Calorimeter covering forward regions},
    first={Liquid Argon Electromagnetic End-Cap (EMEC)},
    text={EMEC}
}

\newglossaryentry{HCal}{
    name=HCal,
    description={Hadronic Calorimeter, an ATLAS calorimeter system designed to measure the energy of hadrons, also known as the Tile Calorimeter},
    first={Hadronic Calorimeter (HCal)},
    text={HCal}
}

\newglossaryentry{MS}{
    name=MS,
    description={Muon Spectrometer, the outermost sub-detector of ATLAS, providing triggering information, identification, and momentum measurements for muons},
    first={Muon Spectrometer (MS)},
    text={MS}
}

\newglossaryentry{SP}{
    name=SP,
    description={Space points, three-dimensional points reconstructed from clustered hits in the ATLAS Inner Detector},
    first={\textit{space points} (SP)},
    text={SP}
}

\newglossaryentry{PV}{
    name=PV,
    description={Primary vertex, a vertex produced by the proton-proton collision in a bunch crossing, typically the hard scatter vertex},
    first={\textit{primary vertex} (PV)},
    text={PV}
}

\newglossaryentry{PFlow}{
    name=PFlow,
    description={Particle Flow, an algorithm used in ATLAS event reconstruction to combine information from the Inner Detector and calorimeters},
    first={Particle Flow (PFlow)},
    text={PFlow}
}

\newglossaryentry{JVT}{
    name=JVT,
    description={Jet Vertex Tagger, a tool in ATLAS that uses associated tracks to discriminate against jets originating from pileup interactions},
    first={Jet Vertex Tagger (JVT)},
    text={JVT}
}

\newglossaryentry{D2obs}{ % Using D2obs as the name for LaTeX compatibility
    name=$\Dtwo$,
    description={The $\Dtwo$ observable ($\sfrac{e_3(\beta)}{e_2(\beta)^3}$), derived from energy correlation functions, commonly used to separate 1-prong and 2-prong jets},
    first=$\Dtwo$,
    text=$\Dtwo$
}

\newglossaryentry{EFPs}{
    name=EFPs,
    description={Energy Flow Polynomials, a set of basis features sensitive to the underlying substructure of different jet types},
    first={Energy Flow Polynomials (EFPs)},
    text={EFPs}
}

\newglossaryentry{DIPS}{
    name=DIPS,
    description={A flavour tagger used by ATLAS based on the deep sets architecture},
    first={DIPS},
    text={DIPS}
}

\newglossaryentry{DL1d}{
    name=DL1d,
    description={A combined flavour tagger used by ATLAS, incorporating the outputs of DIPS with other algorithms},
    first={DL1d},
    text={DL1d}
}

% Deep Learning

\newglossaryentry{ML}{
    name=ML,
    description={Machine Learning, a field that focuses on the development and training of models capable of learning from data},
    first={Machine Learning (ML)},
    text={ML}
}

\newglossaryentry{HEP}{
    name=HEP,
    description={High Energy Physics, the field of physics that studies the fundamental constituents of matter and the interactions between them},
    first={High Energy Physics (HEP)},
    text={HEP}
}

\newglossaryentry{iid}{
    name=i.i.d.,
    description={independent and identically distributed, refers to data samples drawn from the same distribution independently of each other, a common assumption in machine learning training datasets},
    first={independent and identically distributed (i.i.d.)},
    text={i.i.d.}
}

\newglossaryentry{OOD}{
    name=OOD,
    description={Out-of-distribution, refers to data samples that are not drawn from the same distribution as the training data, often used in the context of evaluating model robustness},
    first={out-of-distribution (OOD)},
    text={OOD}
}

\newglossaryentry{MLP}{
    name=MLP,
    description={Multi-Layer Perceptron, arguably the simplest form of a deep neural network, also referred to as a dense network, fully connected network, or sometimes confusingly as a feed-forward neural network (FFN)},
    first={Multi-Layer Perceptron (MLP)},
    text={MLP}
}

\newglossaryentry{FFN}{
    name=FFN,
    description={Feed-Forward Neural Network, a term that can be used to describe any neural network without loops in the computational graph; sometimes used confusingly to refer to a Multi-Layer Perceptron},
    first={feed-forward neural network (FFN)},
    text={FFN}
}

\newglossaryentry{ReLU}{
    name=ReLU,
    description={Rectified Linear Unit, a common non-linear activation function used for hidden layers in modern neural networks, defined as $\max(0, x)$},
    first={ReLU},
    text={ReLU}
}

\newglossaryentry{SGD}{
    name=SGD,
    description={Stochastic Gradient Descent, an iterative optimization algorithm used for training deep neural networks by adjusting parameters in the direction of the negative gradient calculated over a random subset (mini-batch) of the data},
    first={stochastic gradient descent (SGD)},
    text={SGD}
}

\newglossaryentry{Adam}{
    name=Adam,
    description={A common variant of stochastic gradient descent that combines momentum and adaptive learning rates, used for training deep neural networks},
    first={Adam},
    text={Adam}
}

\newglossaryentry{BatchNorm}{
    name=BatchNorm,
    description={Batch Normalization, a type of normalization layer widely used within neural networks to help stabilize training by targeting activations with a mean of zero and unit variance},
    first={batch normalization (BatchNorm)},
    text={BatchNorm}
}

\newglossaryentry{LayerNorm}{
    name=LayerNorm,
    description={Layer Normalization, a type of normalization layer used within neural networks that normalizes activations across the feature dimension},
    first={layer normalization (LayerNorm)},
    text={LayerNorm}
}

\newglossaryentry{PreNorm}{
    name=PreNorm,
    description={Pre-normalization, a setup for residual blocks where the input to the learnable layer is normalized before applying the transformation; helps prevent exploding variance in very deep networks},
    first={pre-normalization (PreNorm)},
    text={PreNorm}
}

% Graph Neural Networks and Variants

\newglossaryentry{GNN}{
    name=GNN,
    description={Graph Neural Network, a type of deep learning model designed to handle graph-structured data, allowing information to flow through the network in a manner that mirrors the graph's organization},
    firstplural={Graph Neural Networks (GNNs)},
    first={Graph Neural Network (GNN)},
    text={GNN}
}

\newglossaryentry{CNN}{
    name=CNN,
    description={Convolutional Neural Network, a class of deep neural networks commonly used for analyzing visual imagery, which applies convolutional operations to exploit spatial structure in the data},
    first={Convolutional Neural Network (CNN)},
    text={CNN},
    firstplural={Convolutional Neural Networks (CNNs)}
}

\newglossaryentry{GCN}{
    name=GCN,
    description={Graph Convolutional Network, a variant of graph neural networks that extends the convolution operation from grid-like data to graph-structured data},
    first={Graph Convolutional Network (GCN)},
    text={GCN}
}

\newglossaryentry{GAT}{
    name=GAT,
    description={Graph Attention Network, a type of graph neural network that incorporates attention mechanisms to weigh the importance of neighboring nodes during message passing},
    first={Graph Attention Network (GAT)},
    text={GAT},
    firstplural={Graph Attention Networks (GATs)}
}

% Transformer Architecture Components

\newglossaryentry{NLP}{
    name=NLP,
    description={Natural Language Processing, a field of machine learning that focuses on the modeling and understanding of human language},
    first={Natural Language Processing (NLP)},
    text={NLP}
}

\newglossaryentry{TE}{
    name=TE,
    description={Transformer Encoder, a variant of the GN Block that operates on a set with no persistent edge or global features, applying self-attention for message passing followed by an MLP for node updates},
    first={Transformer Encoder (TE)},
    text={TE}
}

\newglossaryentry{CV}{
    name=CV,
    description={Computer Vision, a field of machine learning that focuses on enabling computers to interpret and understand visual information},
    first={Computer Vision (CV)},
    text={CV}
}

\newglossaryentry{TD}{
    name=TD,
    description={Transformer Decoder, a component that includes both self-attention and cross-attention mechanisms, allowing one set to be conditioned on another},
    first={Transformer Decoder (TD)},
    text={TD}
}

\newglossaryentry{GNBlock}{
    name=GN Block,
    description={Graph Network Block, a basic building block of a graph neural network, consisting of a message passing step and a node update step},
    first={Graph Network Block (GN Block)},
    text={GN Block}
}

\newglossaryentry{NLNN}{
    name=NLNN,
    description={Non-Local Neural Network, a variant of the GN Block that captures long-range dependencies in data, conceptually similar to the transformer encoder},
    first={Non-Local Neural Network (NLNN)},
    text={NLNN}
}

\newglossaryentry{CA}{
    name=CA,
    description={Cross-Attention, a mechanism that enables conditioning the update of one set on another, creating a one-way bipartite graph between two sets},
    first={Cross-Attention (CA)},
    text={CA}
}

% Attention Mechanisms

\newglossaryentry{SDPA}{
    name=SDPA,
    description={Scaled Dot-Product Attention, the core attention mechanism used in transformers, where messages between nodes are weighted based on the dot product of query and key projections},
    first={Scaled Dot-Product Attention (SDPA)},
    text={SDPA}
}

\newglossaryentry{MHA}{
    name=MHA,
    description={Multi-Headed Attention, an extension of attention where multiple attention operations are performed in parallel with different learned projections, effectively creating a multi-graph on the set},
    first={Multi-Headed Attention (MHA)},
    text={MHA}
}

\newglossaryentry{MHSA}{
    name=MHSA,
    description={Multi-Headed Self-Attention, a specific type of multi-headed attention where the query, key, and value projections all come from the same set of nodes},
    first={Multi-Headed Self-Attention (MHSA)},
    text={MHSA}
}

% Positional Encodings

\newglossaryentry{APE}{
    name=APE,
    description={Absolute Positional Encoding, a technique that integrates sequence positions into node attributes by adding a unique position-dependent tensor to each node before processing by a transformer},
    first={Absolute Positional Encoding (APE)},
    text={APE}
}

\newglossaryentry{RPE}{
    name=RPE,
    description={Relative Positional Encoding, a technique that encodes relative positions between elements in a sequence by biasing the attention matrix, serving as edge information in the graph perspective},
    first={Relative Positional Encoding (RPE)},
    text={RPE}
}

% Other Components

\newglossaryentry{GLU}{
    name=GLU,
    description={Gated Linear Unit, a network layer involving the element-wise multiplication of two linear projections, one of which is activated, commonly used in modern transformer architectures},
    first={Gated Linear Unit (GLU)},
    text={GLU}
}

\newglossaryentry{ParT}{
    name=ParT,
    description={A model developed for jet identification, which uses a transformer architecture to process particle clouds},
    first={Particle Transformer (ParT)},
    text={ParT}
}

\newglossaryentry{ParticleNet}{
    name=ParticleNet,
    description={A graph neural network based jet tagger, which uses a message passing architecture to process particle clouds},
    first={ParticleNet},
    text={ParticleNet}
}

\newglossaryentry{FrechetParticleNetDistance}{
    name=FPND,
    description={Fréchet ParticleNet distance, a metric between two datasets of jets using the Fréchet distance of the penultimate embedding layer of a ParticleNet model},
    first={Fréchet ParticleNet distance (FPND)},
    text={FPND}
}

% Generative Modelling

\newglossaryentry{PGM}{
    name=PGM,
    description={Probabilistic Generative Model, a class of models that either aim to learn the underlying probability distribution of a dataset or at least facilitate means to generate new samples from the learned distribution},
    first={Probabilistic Generative Model (PGM)},
    firstplural={Probabilistic Generative Models (PGMs)},
    text={PGM}
}

\newglossaryentry{NF}{
    name=NF,
    description={Normalizing Flow, a type of generative model that learns a complex distribution by transforming a simple base distribution through a series of invertible transformations},
    first={Normalizing Flow (NF)},
    text={NF},
    firstplural={Normalizing Flows (NFs)}
}

\newglossaryentry{GAN}{
    name=GAN,
    description={Generative Adversarial Network, a type of Probabilistic Generative Model that permits sampling from a density, based on a two-player min-max game between a generator and a discriminator},
    first={Generative Adversarial Network (GAN)},
    firstplural={Generative Adversarial Networks (GANs)},
    text={GAN}
}

\newglossaryentry{VAE}{
    name=VAE,
    description={Variational Autoencoder, an older Probabilistic Generative Model that became popular for generative tasks, unsupervised learning, and variational inference, and is an example of a latent variable model},
    first={Variational Autoencoder (VAE)},
    firstplural={Variational Autoencoders (VAEs)},
    text={VAE}
}

\newglossaryentry{VQVAE}{
    name=VQVAE,
    description={Vector Quantized Variational Autoencoder, a type of VAE that uses vector quantization as the information bottleneck in the latent space,},
    first={Vector Quantized Variational Autoencoder (VQVAE)},
    firstplural={Vector Quantized Variational Autoencoders (VQVAEs)},
    text={VQVAE}
}

\newglossaryentry{ODE}{
    name=ODE,
    description={Ordinary Differential Equation},
    first={ordinary differential equation (ODE)},
    firstplural={ordinary differential equations (ODEs)},
    text={ODE}
}

\newglossaryentry{CNF}{
    name=CNF,
    description={Continuous Normalizing Flow, a type of Normalizing Flow that defines a continuous transformation process using an ordinary differential equation, where each infinitesimal transformation depends on the current state and time},
    first={Continuous Normalizing Flow (CNF)},
    firstplural={Continuous Normalizing Flows (CNFs)},
    text={CNF}
}

\newglossaryentry{SDE}{
    name=SDE,
    description={Stochastic Differential Equation},
    first={stochastic differential equation (SDE)},
    firstplural={stochastic differential equations (SDEs)},
    text={SDE}
}

\newglossaryentry{Sm}{
    name=Sm,
    description={Score Matching, the main objective of a score-based generative model, which involves approximating the score function $\nabla_{\x_t} \log p_t(\x_t)$},
    first={score matching (Sm)},
    text={Sm}
}

\newglossaryentry{DSm}{
    name=DSm,
    description={Denoising Score Matching, an objective for training score-based models that uses the conditional score function $\nabla_{\x_t} \log p_{0t}(\x_t|\x_0)$, which, unlike score function, is tractable},
    first={denoising score matching (DSm)},
    text={DSm}
}

\newglossaryentry{CFM}{
    name=CFM,
    description={Conditional Flow Matching, an alternative framework for constructing Continuous Normalizing Flows that arrives at an equivalent training objective to score-based models},
    first={Conditional Flow Matching (CFM)},
    text={CFM}
}

\newglossaryentry{EDM}{
    name=EDM,
    description={Elucidating the Design Space of Diffusion-Based Generative Models, a diffusion framework where the diffusion schedule is constructed to simplify the form of the probability flow ODE},
    first={EDM},
    text={EDM}
}

\newglossaryentry{EMA}{
    name=EMA,
    description={Exponential Moving Average},
    first={exponential moving average (EMA)},
    text={EMA}
}

\newglossaryentry{CosEnc}{
    name=CosEnc,
    description={Cosine Encoding, a method used for vector embedding scalar variables, involving a layer that outputs a cosine functions with varying increasing frequency},
    first={Cosine Encoding (CosEnc)},
    text={CosEnc}
}

\newglossaryentry{GPU}{
    name=GPU,
    description={Graphics Processing Unit, a specialized processor designed to accelerate graphics rendering and parallel computations},
    first={Graphics Processing Unit (GPU)},
    text={GPU}
}

\newglossaryentry{CPU}{
    name=CPU,
    description={Central Processing Unit, the primary component of a computer that performs most of the processing inside a computer},
    first={Central Processing Unit (CPU)},
    text={CPU}
}

\newglossaryentry{SVD}{
    name=SVD,
    description={Singular Value Decomposition method, a form of matrix factorization},
    first={Singular Value Decomposition (SVD)},
    text={SVD}
}

\newglossaryentry{FSR}{
    name=FSR,
    description={Final State Radiation, the emission of additional particles from the final state of a particle interaction},
    first={Final State Radiation (FSR)},
    text={FSR}
}

\newglossaryentry{ISR}{
    name=ISR,
    description={Initial State Radiation, the emission of additional particles from the initial state of a particle interaction},
    first={Initial State Radiation (ISR)},
    text={ISR}
}

\newglossaryentry{NuWeight}{
    name=\vweight,
    description={An analytical technique for reconstructing neutrino kinematics, typically in multi-neutrino events, based on a grid search over rapidities and using mass constraints},
    first={\vweight},
    text={\vweight}
}

\newglossaryentry{NuEllipse}{
    name=\ellipse,
    description={A geometric approach for reconstructing neutrino kinematics using mass constraints, particularly in processes where top quarks decay leptonically},
    first={\ellipse},
    text={\ellipse}
}

\newglossaryentry{MC}{
    name=MC,
    description={Monte Carlo, a method used in simulations to model complex systems and processes by generating random samples from a probability distribution},
    first={Monte Carlo (MC)},
    text={MC}
}

\newglossaryentry{JetNet}{
    name=JetNet,
    description={An open dataset of simulated \pythia jets,},
    first={JetNet},
    text={JetNet}
}

\newglossaryentry{JetClass}{
    name=JetClass,
    description={An open dataset of simulated \pythia+\delphes jets},
    first={JetClass},
    text={JetClass}
}

\newglossaryentry{tauij}{
    name=$\tau_{ij}$,
    description={A substructure variable, specifically the ratio of the $i$ and $j$ N-subjettiness values, used to distinguish jets based on their prong structure},
    first={$\tau_{ij}$},
    text={$\tau_{ij}$}
}

\newglossaryentry{CWoLa}{
    name=CWoLa,
    description={Classification Without Labels, a semi-supervised learning method for training classifiers using datasets with mixed levels of signal and background},
    first={Classification Without Labels (CWoLa)},
    text={CWoLa}
}

\newglossaryentry{Cathode}{
    name=Cathode,
    description={A anomaly detection method where a conditional generative model is used to construct a reference background template which is then used to train a classifier using CWoLa},
    first={Cathode},
    text={Cathode}
}

\newglossaryentry{Drapes}{
    name=Drapes,
    description={A class of models which use conditional diffusion to generate background templates for anomaly detection, based on the \pcdroid model},
    first={\drapes},
    text={\drapes}
}

\newglossaryentry{DrapesHLV}{
    name=Drapes HLV,
    description={A variant of the Drapes model that generates only high-level features for anomaly detection templates},
    first={\drapes HLV},
    text={\drapes HLV}
}

\newglossaryentry{DrapesLLV}{
    name=Drapes LLV,
    description={A variant of the Drapes model that generates low-level features (di-jet system as sets of constituents) for anomaly detection templates},
    first={\drapes LLV},
    text={\drapes LLV}
}

\newglossaryentry{SignificanceImprovement}{
    name=SIC,
    description={Significance Improvement, A metric for anomaly detection classifiers, defined as the number of signal events divided by the square root of the number of background events passing the classifier filter},
    first={significance improvement (SIC)},
    text={SIC}
}

\newglossaryentry{PIPPIN}{
    name=PIPPIN,
    description={Particles Into Particles with Permutation Invariant Networks, a novel method that frames full forward simulation as a set-to-set generative task, directly generating reconstructed objects from partons},
    first={PIPPIN (Particles Into Particles with Permutation Invariant Networks)},
    text={PIPPIN}
}

\newglossaryentry{ConsistencyDistillation}{
    name=CD,
    description={Consistency distillation, a process where a teacher network trains a student network to solve the reverse diffusion process in fewer steps, potentially enabling one-step generation},
    first={consistency distillation (CD)},
    text={CD}
}

\newglossaryentry{NFE}{
    name=NFE,
    description={Neural Function Evaluations, a measure of the number of network passes required during generation with diffusion models, directly impacting generation time},
    first={neural function evaluations (NFE)},
    text={NFE}
}

\newglossaryentry{Pythia}{
    name=Pythia,
    description={A general-purpose Monte Carlo event generator for simulating high-energy particle collisions, including parton showering, hadronization, and multi-parton interactions},
    first={\pythia},
    text={\pythia}
}

\newglossaryentry{Sherpa}{
    name=Sherpa,
    description={A flexible event generator designed for precision QCD calculations, incorporating matrix elements, parton showers, and hadronization to model complex collider events},
    first={\sherpa},
    text={\sherpa}
}

\newglossaryentry{Herwig}{
    name=Herwig,
    description={A comprehensive event generator focusing on accurate modeling of parton showers and non-perturbative QCD effects, with strong capabilities in hadronization and BSM physics},
    first={\herwig},
    text={\herwig}
}

\newglossaryentry{MadGraph}{
    name=MadGraph,
    description={A software package for simulating high-energy particle collisions, focusing on matrix element generation and event generation for various processes},
    first={\madgraph},
    text={\madgraph}
}

\newglossaryentry{akT}{
    name=anti-$k_T$,
    description={A jet clustering algorithm used to group particles into jets, prioritizing the angular separation of particles},
    first={anti-$k_T$},
    text={anti-$k_T$}
}

\newglossaryentry{GEANT4}{
    name=GEANT4,
    description={A state-of-the-art tool for high-precision detector simulations in particle physics},
    first={\geant},
    text={\geant}
}

\newglossaryentry{Delphes}{
    name=Delphes,
    description={A fast parametric simulation tool used as an efficient alternative to computationally expensive detector simulations when resources are limited},
    first={\delphes},
    text={\delphes}
}

\newglossaryentry{FM}{
    name=FM,
    description={Foundation Model, a deep learning model that serves as the basis for further development or research, often produced using unsupervised or self-supervised learning techniques},
    first={Foundation Model (FM)},
    text={FM}
}

\newglossaryentry{BERT}{
    name=BERT,
    description={Bidirectional Encoder Representations from Transformers, a transformer-based model pre-trained on large text corpora, widely used in NLP tasks},
    first={Bidirectional Encoder Representations from Transformers (BERT)},
    text={BERT}
}

\newglossaryentry{GPT}{
    name=GPT,
    description={Generative Pre-trained Transformer, a transformer-based model pre-trained on large text corpora, widely used in NLP tasks},
    first={Generative Pre-trained Transformer (GPT)},
    text={GPT}
}

\newglossaryentry{BEiT}{
    name=BEiT,
    description={Bidirectional Encoder representation from Image Transformers, a transformer-based model pre-trained on large image corpora, widely used in computer vision tasks},
    first={Bidirectional Encoder representation from Image Transformers (BEiT)},
    text={BEiT}
}

\newglossaryentry{MAE}{
    name=MAE,
    description={Masked Autoencoder, a transformer-based model pre-trained on large image corpora, widely used in computer vision tasks},
    first={Masked Autoencoder (MAE)},
    text={MAE}
}



